% ======================================================================
% SECTION 1: INTRODUCTION
% File: sections/introduction.tex
% ======================================================================

[PROCESSING INSTRUCTION: This introduction begins on the title page directly below the abstract and
keywords. It should be approximately 1.5 pages. The first 10-12 lines appear on the title page. The
remaining content continues after the table of contents on page 3. Generate the introduction using
content from the following source files:

SOURCE FILES TO PROCESS:
- sponsor/input\_files/sponsor\_01\_executive\_strategy.md (executive vision, portfolio strategy)
- sponsor/input\_files/org\_01\_executive\_regulatory.md (sponsor accountability under 21 CFR 312.50/312.52)
- national-platform/new\_paper/final\_paper/sections/introduction.tex (platform context)
- national-platform/new\_paper/final\_paper/sections/executive\_summary.tex (platform overview)

CONTENT REQUIREMENTS:
1. Open with the challenge: traditional oncology trial sponsors employ hundreds of staff across
   12+ functional areas (clinical ops, safety, regulatory, quality, supply chain, etc.) costing
   significant overhead and introducing human bottlenecks in trial execution.
2. Introduce the paradigm shift: Physical AI systems (surgical robots, digital twins, federated
   learning) demand a new sponsor operating model that matches their autonomous capabilities.
3. State the thesis: this paper presents a fully autonomous AI-native sponsor operating system
   that replaces each human sponsor function with a software agent or robotic subsystem.
4. Reference the regulatory foundation: adapted 21 CFR Part 312, 21 CFR Part 50, and ICH E6(R3)
   from prior work in this repository (cite references from references.bib).
5. Describe the multi-agent architecture overview: 12 functional agents organized into governance,
   execution, site/robotics, and trust layers.
6. Clarify what the system is: a company-shaped trial execution application, not merely an EDC,
   CTMS, robot controller, protocol copilot, safety dashboard, or site orchestration tool.
7. Clarify what it is not: not a replacement for the legal sponsor of record, but an AI-native
   operating system wrapped by a human/legal sponsor of record.
8. Reference the existing platform infrastructure: MCP servers, federated learning, PSL/USL scoring,
   site documentation package, and patient journey orchestration.
9. End with a roadmap of the paper sections.

PYTHON SCRIPT REQUIREMENTS (to be generated by Claude Code Opus 4.6 in the future):
- File: sponsor/template/scripts/generate\_introduction\_metrics.py
- Purpose: Parse sponsor/input\_files/sponsor\_01\_executive\_strategy.md and
  org\_01\_executive\_regulatory.md to extract key regulatory dates, cost figures,
  and staffing numbers. Output a JSON summary used to populate introduction statistics.
- Must handle markdown parsing, regex extraction of dollar amounts and dates.
- Process with Claude Code Opus 4.6 (1M token context).

FORMATTING NOTES:
- First paragraph should be compelling and self-contained for the title page portion.
- Use \S for section symbol references (not SS or double-S).
- Ensure no text runs off the right margin. Use proper hyphenation.]

The integration of Physical AI systems into oncology clinical trials has established a new
paradigm for cancer treatment evaluation, delivery, and monitoring. Prior work in this
repository has developed the regulatory foundations (adapted 21 CFR Part 312, 21 CFR Part 50,
and ICH E6(R3)), quantitative readiness standards (PSL and USL scoring), national MCP server
infrastructure, federated learning frameworks, site establishment documentation, and a
complete single-patient journey orchestration. These developments address the trial execution
layer but leave the sponsor layer, the organizational entity that designs, funds, governs,
and takes accountability for the trial, operating under a traditional human-staffed model.

[PROCESSING INSTRUCTION: Continue the introduction here after the table of contents. Expand
the above opening paragraph and add 4-5 additional paragraphs covering the points listed above.
Reference the staffing scale table from org\_02\_operating\_model\_roles.md showing 10 functional
areas. Reference the 7 governance forums from org\_02. Reference the 7 decision gates from
sponsor\_07\_timelines\_kpis.md. Include budget benchmarks from sponsor\_07: Phase I \$14.2M,
Phase II \$25.5M, Phase III \$65.8M (Tufts CSDD 2022). Include daily delay costs: Phase III
\$55,716/day. End with: ``The remainder of this paper is organized as follows. \S 2 presents...''
listing all 18 sections. Process with Claude Code Opus 4.6 (1M token context).]
